\chapter*{CONCLUSIONES}
\phantomsection
\addcontentsline{toc}{chapter}{CONCLUSIONES}

El presente trabajo tuvo como objetivo diseñar e implementar un prototipo de robot móvil diferencial capaz de generar mapas 2D, ejecutar navegación autónoma y realizar exploración basada en fronteras mediante \gls{ROS2} Foxy. El desarrollo integró diseño mecánico, electrónica de potencia, firmware de bajo nivel, gemelo digital en Gazebo Classic, mapeo con \texttt{slam\_toolbox}, navegación con Nav2 y exploración autónoma.

En relación con la plataforma física, se construyó un robot diferencial de bajo
costo con la configuración descrita en el Capítulo~\ref{ch:implementacion}, en
la que el \gls{Arduino} ejecuta la lectura de encoders y el control de cada
rueda, la Raspberry Pi ejecuta la interfaz de hardware, el controlador
diferencial y la adquisición del \gls{LIDAR}, y la estación de trabajo ejecuta
SLAM, navegación, exploración y supervisión.

El gemelo digital desarrollado en Gazebo Classic permitió verificar la geometría,
la cinemática, la cadena de transformaciones y la interfaz de control antes de
operar el robot físico. En simulación, la desviación media respecto de la
referencia nominal de 2~m fue 3.39\%, calculada con la propia odometría porque
la pose verdadera de Gazebo Classic no quedó registrada durante los ensayos. Ese valor
no separa el error odométrico del error de seguimiento. En hardware real, el
error odométrico medio fue 3.63\% y el error físico de seguimiento fue 7.18\%.
La selección del backend y su trazabilidad se documentan en el
\ref{anx:gemelo_digital}.

Respecto del control de bajo nivel, el firmware ejecuta de forma independiente
para cada rueda la ley de control heredada del firmware original, que
corresponde a un controlador PI en forma incremental con la acción
proporcional sobre la medición, más un término integral adicional, y cuyas
ganancias no cumplen la función que indica su nombre. La
Fig.~\ref{fig:respuesta_escalon_final} corresponde a una adquisición
experimental en lazo cerrado, y los indicadores de la
Tabla~\ref{tab:pid_performance} se calcularon sobre esa misma respuesta medida.
Ambas ruedas alcanzan la referencia sin sobreimpulso apreciable. Los estados internos, la temporalidad y las aproximaciones
numéricas se documentan en los anexos de implementación y sintonización.

El mapeo con \texttt{slam\_toolbox} permitió construir mapas 2D en simulación y hardware real. Los indicadores de cobertura libre y de área mapeada se cuantificaron sobre hardware real durante las iteraciones de exploración.
La navegación con Nav2 quedó integrada mediante planificación global,
control local, mapas de costo y acciones de recuperación. La exploración por fronteras se ejecutó en cinco iteraciones válidas, que
terminaron automáticamente al no quedar fronteras y sin intervención manual, con
un índice medio de cobertura libre de 62.2\% ($\sigma=12.9$). De las 60 metas
generadas, 54 quedaron abortadas y una se completó. En 53 de esos abortos la
meta fue reemplazada por el propio explorador, que renueva su destino cada
3.3~s, y Nav2 informa ese reemplazo como aborto. El detalle metodológico se conserva en el
\ref{anx:resultados_metodologia}.
El cumplimiento de los objetivos específicos se resume de la siguiente forma:

\begin{description}
\item[\autoref{oe1} -- Construcción y gemelo digital:] \textit{Cumplido.} Se diseñó y construyó la plataforma diferencial física y se verificó su gemelo digital con el backend de \texttt{ros2\_control} ejecutado en Gazebo Classic.

\item[\autoref{oe2} -- Control de bajo nivel:]
\textit{Cumplido parcialmente.} Se implementaron la lectura de encoders,
la estimación de los estados de rueda y el control de velocidad en el firmware;
la posición y velocidad de rueda se exportan mediante
\texttt{DiffDriveArduino}, y la odometría se integra en
\texttt{diff\_drive\_controller}. Sin embargo, la ley heredada de
\texttt{ROSArduinoBridge} opera como un controlador PI incremental sin acción
derivativa efectiva y se evaluó en un único punto de operación, como se detalla
en la Sec.~\ref{sec:motorcontrol_pid} y en la Sec.~\ref{sec:resultados_pid}.

\item[\autoref{oe3} -- Localización, mapeo y navegación:]
\textit{Cumplido parcialmente.} El mapeo y la localización durante la
construcción del mapa se implementaron y verificaron funcionalmente en hardware
real, junto con la caracterización cuantitativa de cobertura y área libre
mapeada. Nav2 se integró y condujo al robot hacia las metas del explorador en
las cinco ejecuciones válidas, sin abortos atribuibles a fallos de navegación en
el registro. Sin embargo, como 53 de los 54 abortos corresponden a reemplazos de
meta hechos por el explorador, el resultado de las metas no permite medir la
tasa de éxito de la navegación. El carácter parcial proviene de que no se obtuvo
un indicador cuantitativo del desempeño de la navegación.

\item[\autoref{oe4} -- Exploración autónoma:] \textit{Cumplido.} La estrategia
de exploración por fronteras se integró y ejecutó sobre el robot real, y generó
60 metas hacia regiones no exploradas sin intervención manual. Cinco de las seis
ejecuciones terminaron de forma automática por el criterio del propio
explorador, al no quedar fronteras en el mapa de costos, y produjeron mapas
válidos con un índice medio de cobertura libre de 62.2\% y un área libre mapeada
media de 3.50~m$^2$. En la ejecución restante el robot no se desplazó y la causa
no se determinó. La exploración quedó verificada funcionalmente y caracterizada,
sin una referencia geométrica externa para validar el mapa. Su principal
limitación es que el explorador agrega a la lista negra las metas que él mismo
reemplaza, efecto de la integración con Nav2 que no impidió terminar la
exploración.

Respecto de la pregunta de investigación planteada en la Introducción, los
resultados muestran que fue posible construir un prototipo diferencial capaz de
generar mapas 2D y ejecutar un proceso de exploración autónoma en el entorno
ensayado. Cinco de las seis ejecuciones experimentales produjeron mapas válidos,
lo que permite verificar funcionalmente la integración entre exploración,
navegación y mapeo.

El resultado de las metas individuales no contradice esa conclusión, porque 53
de los 54 abortos corresponden a reemplazos de meta hechos por el propio
explorador. Sí deja pendiente medir cuantitativamente el desempeño de la
navegación, lo que requiere un protocolo con metas fijas.

\section*{Limitaciones}
\phantomsection
\addcontentsline{toc}{section}{Limitaciones}

\begin{itemize}
\item \textbf{Deriva odométrica:} la odometría basada solo en encoders acumula error, especialmente en giros y trayectorias largas.
\item \textbf{Espacio de pruebas reducido:} el recinto disponible limita la ejecución de recorridos extensos y condiciona la geometría de los mapas de costo.
\item \textbf{Registro de navegación:} los contadores automáticos de la sesión no identifican metas únicas, por lo que el resultado de las metas de exploración se reconstruyó a posteriori desde el flujo de estado de la acción. No se ejecutó un protocolo de navegación dirigida con metas fijas, por lo que no se dispone de tasa de éxito ni de error de posición final para ese caso.
\item \textbf{Métricas PID:} los indicadores de la Tabla~\ref{tab:pid_performance} provienen de la captura experimental en hardware real, mientras que los sobreimpulsos de 10.7\% a 15.6\% provienen de una simulación ilustrativa sobre un modelo FOPDT reidentificado. No deben interpretarse como una única prueba.
\item \textbf{Fidelidad de la simulación:} Gazebo Classic no reproduce completamente fricción, zona muerta del driver, ruido sensorial ni latencias reales.
\item \textbf{Distribución en fin de vida:} \gls{ROS2} Foxy ya no recibe soporte oficial.
\item \textbf{Validez externa:} los resultados provienen de un único recinto, cuya área libre mapeada media fue de 3.50~m$^2$, y de cinco iteraciones, por lo que no se generalizan a entornos mayores o de geometría distinta. La superficie del recinto no se midió de forma independiente del mapa.
\item \textbf{Ley de control de bajo nivel:} el firmware ejecuta la ley acumulativa heredada de \texttt{ROSArduinoBridge}, en la que $K_p$ actúa como ganancia integral, $K_d$ como ganancia proporcional sobre la medición y $K_i$ como segunda integración (Ec.~\eqref{eq:pid_equivalente}). No se realizó un análisis de estabilidad de esta ley ni una comparación con un PID incremental estándar.
\item \textbf{Particularidades de implementación conservadas:} la lectura no atómica de los encoders, la posición global del PSO copiada desde la posición actual y sin semilla fija, y el costo de fronteras no homogéneo en unidades se documentan tal como se ejecutaron. Su efecto sobre los resultados no se cuantificó.
\item \textbf{Integración entre explorador y Nav2:} el explorador envía una meta nueva sin cancelar la anterior, y Nav2 informa la meta reemplazada como abortada. Por eso el resultado de las metas no mide la tasa de éxito de la navegación, y el explorador agrega a su lista negra fronteras que no fallaron. El efecto sobre la duración de la exploración no se cuantificó.
\end{itemize}

\section*{Trabajo futuro}
\phantomsection
\addcontentsline{toc}{section}{Trabajo futuro}

\begin{enumerate}
\item Ejecutar el protocolo cuantitativo de navegación en un entorno de mayor superficie, registrando cada meta y su resultado de forma unívoca.
\item Integrar una \gls{IMU} mediante el paquete de fusión sensorial \texttt{robot\_localization} para reducir la deriva angular.
\item Migrar el sistema a una distribución mantenida de \gls{ROS2} y portar en el mismo paso el gemelo digital a una versión vigente de Gazebo, dado que Gazebo Classic ya no recibe soporte.
\item Evaluar la navegación en espacios estrechos con una configuración acorde a la geometría del entorno.
\item Registrar por separado los reemplazos de meta y los fallos de Nav2, cancelando la meta en curso antes de enviar la siguiente o conservando el registro de \texttt{bt\_navigator}, y evaluar el efecto de \texttt{planner\_frequency} sobre la duración de la exploración.
\item Extender la evaluación del control de velocidad a otras referencias y
compararlo con un PID incremental estándar, incorporando un estado integral en
punto flotante y lectura atómica de encoders.
\item Extender el sistema a la localización sobre un mapa almacenado mediante AMCL, comparando la pose estimada con posiciones marcadas físicamente en el recinto.
\item Repetir el protocolo angular con un giro controlado a un ángulo fijo y verificar la velocidad angular efectiva registrada en \texttt{/odom}.
\item Registrar la pose verdadera de Gazebo Classic en los ensayos simulados para separar el error odométrico del error de seguimiento en simulación.
\end{enumerate}